It is easy to see that the space $\widetilde L^p(\Omega)$ is not complete (see Problem~\ref{ex:continuous-lp-not-complete}). We therefore make the following definition.

\begin{Definition}
\label{def:lp-space-completion}
Let $1\leq p<\infty$. Then $L^p(\Omega)$ is the completion of $\widetilde L^p(\Omega)$.
\end{Definition}

\begin{Remark}
\label{rem:lp-completion-defines-distributions}
If a sequence $f_n$ is Cauchy in $\widetilde L^p(\Omega)$, then it follows from Hölder\textquotesingle s inequality that $\int_\Omega f_n\phi\,\mathrm{d}\bm{x}$ converges for every test function $\phi$. Hence $f_n$ converges in the sense of distributions, and by Theorem~\ref{thm:distribution-sequential-completeness} every element of $L^p(\Omega)$ defines a distribution.
\end{Remark}

We need to show that different elements of $L^p(\Omega)$ correspond to different distributions.

\begin{Lemma}
\label{lem:lp-completion-faithful-distribution-inclusion}
Let $f_n$ be a Cauchy sequence in $\widetilde L^p(\Omega)$ $(1\leq p<\infty)$ and assume that $\int_\Omega f_n\phi\,\mathrm{d}\bm{x}\to0$ for every $\phi\in\mathcal{D}(\Omega)$. Then $\|f_n\|_p\to0$.
\end{Lemma}

\begin{Proof}
Assume the contrary, i.e., let $\lim_{n\to\infty}\|f_n\|_p=M>0$. Choose $N$ large enough so that $3M/2>\|f_n\|_p>M/2$ and $\|f_n-f_m\|_p<\epsilon$ for $n,m\geq N$. Here $\epsilon$ is chosen small. By Problem~\ref{ex:test-functions-dense-lp}, $\mathcal{D}(\Omega)$ is dense in $L^p(\Omega)$. Hence we can choose $\psi\in\mathcal{D}(\Omega)$ with $\|f_N-\psi\|_p<\epsilon$; more specifically, we can choose $\psi$ to be the product of a polynomial and a cutoff function. Now let
\begin{equation}\label{eq:lp-faithful-distribution-dual-test-function}
g=\begin{cases}|\psi|^{p-1}\operatorname{sgn}\psi&p>1\\\operatorname{sgn}\psi&p=1.\end{cases}
\end{equation}
Thus $g$ is a piecewise continuous function of compact support. If $p>1$, we have $g\in L^q(\Omega)$, and we can choose $\phi\in\mathcal{D}(\Omega)$ with $\|g-\phi\|_q<\epsilon$; here $q=p/(p-1)$. We now estimate, for $n\geq N$,
\begin{equation}\label{eq:lp-faithful-distribution-positive-lower-bound}
\begin{aligned}
\left|\int_\Omega f_n\phi\,\mathrm{d}\bm{x}\right|
&\geq\int_\Omega\psi g\,\mathrm{d}\bm{x}\\
&\quad-\left|\int_\Omega\psi(\phi-g)+(f_n-f_N)\phi+(f_N-\psi)\phi\,\mathrm{d}\bm{x}\right|\\
&\geq\int_\Omega|\psi|^p\,\mathrm{d}\bm{x}-\|\psi\|_p\|\phi-g\|_q\\
&\quad+\|f_n-f_N\|_p\|\phi\|_q+\|f_N-\psi\|_p\|\phi\|_q\\
&\geq\|\psi\|_p^p-O(\epsilon)\\
&\geq(M/2-\epsilon)^p-O(\epsilon).
\end{aligned}
\end{equation}
Since $\epsilon$ can be chosen arbitrarily small for large $N$, this contradicts the fact that the left-hand side of \eqref{eq:lp-faithful-distribution-positive-lower-bound} should tend to zero as $n\to\infty$.

If $p=1$, we choose $\phi$ such that $\max|\phi|\leq1$ and such that $\phi$ is uniformly close to $g$ except in a small neighborhood of those surfaces where $g$ is discontinuous. We can then make $\int\psi(\phi-g)\,\mathrm{d}\bm{x}$ as small as we wish. The rest goes as above, with the elementary inequality $|\int fg\,\mathrm{d}\bm{x}|\leq\max|g|\|f\|_1$ in place of Hölder\textquotesingle s inequality.
\end{Proof}

\begin{Remark}
\label{rem:lp-point-values-and-measure-zero}
Consider the sequence $f_n(x)=\exp(-nx^2)$. It is easy to see that $f_n\to0$ in $\widetilde L^p(\mathbb{R})$ for $1\leq p<\infty$, and thus the sequence can be identified with the constant function $0\in L^p(\mathbb{R})$. However, note that $f_n(0)=1$ for every $n$. This example should convince the reader that we cannot regard elements of $L^p$ as functions with point values. Fortunately, using the results of measure theory, it can be shown that one can ``almost'' do so. In measure theory, elements of $L^p$ are defined as equivalence classes of functions which differ only on ``sets of measure zero'' (for instance, at a finite number of points). More precisely, $L^p(\Omega)$ is the set of all (equivalence classes of) Lebesgue measurable functions $f$ such that $|f|^p$ is integrable. This definition is equivalent to ours. To show this equivalence, it suffices to prove that continuous bounded functions are dense in the space $L^p$ (as defined in integration theory). Since we do not assume familiarity with Lebesgue integration, we shall not pursue this point further.
\end{Remark}

On the basis of our definition of $L^p(\Omega)$, we can define ``integrals'' as limits. For example, if $f\in L^1(\Omega)$, then there is a sequence $f_n\in\widetilde L^1(\Omega)$ such that $f_n\to f$, and we can define
\begin{equation}\label{eq:lp-completion-integral-definition}
\int_\Omega f\,\mathrm{d}\bm{x}=\lim_{n\to\infty}\int_\Omega f_n\,\mathrm{d}\bm{x}.
\end{equation}
Also, if $f\in L^p(\Omega)$ and $g\in L^q(\Omega)$, where $p,q\in(1,\infty)$ and $p^{-1}+q^{-1}=1$, we can choose sequences $f_n\in\widetilde L^p(\Omega)$ and $g_n\in\widetilde L^q(\Omega)$ such that $f_n\to f$ and $g_n\to g$. Hölder\textquotesingle s inequality implies that $f_ng_n$ converges in $L^1(\Omega)$, and it is natural to call the limit $fg$. With this interpretation Hölder\textquotesingle s inequality holds for all $f\in L^p(\Omega)$ and $g\in L^q(\Omega)$. These examples are typical of the distributional approach. Properties of functions in $L^p$ and similar spaces are proved by first considering continuous (or even smoother) functions and then considering the limit of an appropriate sequence.

\Needspace{4\baselineskip}
\begin{Example}
\label{ex:banach-direct-product}
Let $X_1$ and $X_2$ be Banach spaces. Then $X_1\times X_2$ is a Banach space with the norm $\|(x_1,x_2)\|=\|x_1\|+\|x_2\|$. In this situation, we shall sometimes identify $X_1$ with $X_1\times\{0\}$ and $X_2$ with $\{0\}\times X_2$. We then write $x_1+x_2$ for $(x_1,x_2)$.
\end{Example}
